% ============================================================================
% ВЕРСИЯ-ЗАГЛУШКА — типографируется, пока \vtwofalse в preamble.tex
%
% Это плейсхолдер, появляющийся в скомпилированной статье, пока лестница
% композиции ещё выполняется. Он формулирует дизайн и предрегистрированное
% предсказание и намеренно не содержит НИ ОДНОГО измеренного числа.
%
% ЧТОБЫ АКТИВИРОВАТЬ РЕАЛЬНЫЙ РАЗДЕЛ: установить \vtwotrue в preamble.tex.
% Реальный текст живёт в 09_results_composition.tex.
% ============================================================================

\section{Композиция состояния: что должно содержать состояние рантайма}
\label{sec:composition}

Разделы~\ref{sec:capability}--\ref{sec:economics} устанавливают, что ущерб
происходит от отброшенной истории и что агент платит за него, заново
получая информацию через свои инструменты. Этот диагноз подразумевает
исправление: если бы состояние несло факты, которые агент постоянно заново
выводит, повторное приобретение не понадобилось бы.

Два свойства измеренного отказа мотивируют дизайн. Неуспешное состояние
пусто ровно от тех фактов, которые заново выводятся (ноль точных версий,
ноль текста ошибок, ноль команд (\S\ref{sec:mechanism})), а собственная
таксономия ошибок оригинала приписывает \SI{68}{\percent} отказов моделей с
открытыми весами перезаписи моделью собственного состояния\origin. Поле,
которое пишет рантайм, не может быть перезаписано моделью, а факт, который
выводит рантайм, не нужно перефразировать по выходному тарифу.

\paragraph{Фактор.} Композиция состояния варьируется как единственный
упорядоченный фактор, по одному компоненту на ступень, всё остальное
заморожено ($\tailk{=}0$, одна схема, без эскалации):

\begin{center}
\begin{tabular}{@{}llp{7.2cm}@{}}
\toprule
Ступень & Добавляет & Вопрос, который изолирует \\
\midrule
v1 & имена файлов, свежесть, счётчики изменений & достаточно ли состояния в форме указателей? \\
v2 & $+$ содержимое пронаблюдённых файлов        & выигрывает ли перенос байтов у переноса координат? \\
v3 & $+$ команды с кодами выхода         & помогает ли знание того, что уже выполнено? \\
v4 & $+$ диагностика с её выводом    & стоит ли своих токенов сохранённый результат компилятора? \\
v5 & $+$ факты об окружении                & стоят ли устойчивые проверки кэширования? \\
v6 & $+$ разобранный исход теста              & добавляет ли что-либо разобранный вердикт сверх сырого вывода? \\
\bottomrule
\end{tabular}
\end{center}

Переключатель «всё или ничего» не может ответить ни на один из этих
вопросов: выигрыш можно было бы приписать шести решениям дизайна сразу, а
проигрыш скрыл бы, какой компонент работал. Лестница принудительна, а не
просто задумана: тесты падают, если ступень добавляет больше одного
компонента, если две ступени отрисовывают идентичные блоки, или если
реализованная ступень запуска расходится со ступенью, с которой он был
запущен\code.

\paragraph{Предрегистрированное предсказание.} Если перенос содержимого —
правильный рычаг, то относительно ступени только-указателей ступени с
содержимым должны сократить долю избыточного осмотра с измеренных
\num{0.58} в сторону контрольных \num{0.20}, \emph{и} это сокращение должно
быть каналом, через который приходит любой выигрыш в доле решённых задач.
Если доля решённых задач движется, а избыточность нет, выигрыш пришёл
откуда-то ещё, и объяснение в этой статье было бы неверным.

\pending{Лестница выполняется. Испытание~1 по всему корпусу завершено, и
его направление согласуется с предсказанием, но одно испытание не
приводится как результат в этой работе по причине, данной в
\S\ref{sec:pitfalls}: эффект, значимый при двух испытаниях в этом проекте,
уже умирал на третьем. Когда набор испытаний закроется, установить
\texttt{\textbackslash vtwotrue} в \texttt{preamble.tex}, и раздел~09
заменит эту заглушку.}
